\section{Tổng hợp lipid -- carbohydrate -- protein}

\subsection{Hệ thống hóa ba nhóm hợp chất}

Lipid, carbohydrate và protein đều là những nhóm hợp chất hữu cơ quan trọng trong cơ thể sống và thực phẩm. Có thể phân biệt chúng dựa trên đặc điểm cấu tạo, tính chất và phản ứng đặc trưng.

\begin{center}
\begin{tabularx}{0.98\textwidth}{>{\bfseries}p{2.15cm}YY}
\toprule
Nhóm chất & Đặc điểm cấu tạo/nguồn thường gặp & Tính chất và dấu hiệu nhận biết tiêu biểu \\
\midrule
Lipid (chất béo) &
Chất béo đơn giản là triester của glycerol và các acid béo; có trong dầu thực vật, mỡ động vật. &
Thường không tan trong nước. Chất béo phản ứng với kiềm khi đun nóng tạo glycerol và muối của acid béo (xà phòng). \\
Carbohydrate &
Gồm các chất như glucose, saccharose, tinh bột và cellulose; có nhiều trong quả chín, mía, gạo, ngô, khoai, gỗ và rơm rạ. &
Glucose cho phản ứng tráng bạc; saccharose bị thủy phân; hồ tinh bột tạo màu xanh tím với iodine; tinh bột và cellulose có thể thủy phân tạo glucose trong điều kiện thích hợp. \\
Protein &
Được tạo thành từ nhiều amino acid liên kết với nhau bằng liên kết peptide; có trong thịt, cá, trứng, sữa, đậu và hạt. &
Có thể bị thủy phân thành amino acid; bị biến tính/đông tụ dưới tác dụng của nhiệt độ hoặc acid/base thích hợp. \\
\bottomrule
\end{tabularx}
\end{center}

Không nên xem mọi carbohydrate có cùng một phép thử: iodine nhận biết tinh bột, còn phản ứng tráng bạc được dùng để nhận biết glucose trong phạm vi các mẫu đã biết. Saccharose không cho phản ứng tráng bạc khi thử trực tiếp.

\subsection{Lựa chọn phép thử và thiết kế trình tự nhận biết}

Khi nhận biết các mẫu chưa rõ tên, cần:
\begin{enumerate}
\item Xác định những chất có thể có trong bộ mẫu và chọn tính chất khác biệt nhất.
\item Chia mỗi mẫu thành nhiều phần nhỏ riêng biệt; không dùng lại phần mẫu đã trộn với thuốc thử.
\item Chọn thuốc thử theo mục tiêu, ghi rõ hiện tượng dương tính và hiện tượng không xuất hiện.
\item Dùng thêm một phép thử xác nhận khi phép thử đầu chỉ cho kết luận sơ bộ.
\end{enumerate}

\begin{center}
\begin{tabularx}{0.98\textwidth}{>{\bfseries}p{2.35cm}YY}
\toprule
Mục tiêu & Phép thử gợi ý & Hiện tượng cần quan sát \\
\midrule
Tinh bột &
Nhỏ dung dịch iodine vào hồ tinh bột hoặc mẫu thử phù hợp. &
Mẫu chứa tinh bột chuyển xanh tím. \\
Glucose &
Cho mẫu vào thuốc thử Tollens rồi đun nóng nhẹ theo hướng dẫn của giáo viên. &
Xuất hiện bạc kim loại bám trên thành ống nghiệm. \\
Protein &
Đun nóng mẫu protein như lòng trắng trứng hoặc thêm một lượng nhỏ acid loãng vào mẫu phù hợp. &
Protein có thể biến tính, đông tụ hoặc kết tủa. \\
Chất béo &
Thử khả năng không tan trong nước; nếu cần phân biệt triglyceride với dầu bôi trơn, đun nóng với dung dịch kiềm. &
Dầu ăn tạo lớp riêng với nước và tham gia phản ứng xà phòng hóa; dầu bôi trơn có bản chất hydrocarbon không tạo xà phòng theo phản ứng này. \\
\bottomrule
\end{tabularx}
\end{center}

\begin{ghinho}{Lưu ý khi kết luận}
\begin{itemize}
\item Một hiện tượng âm tính không phải lúc nào cũng chứng minh chất hoàn toàn không có mặt; mẫu có thể quá loãng, bị che màu hoặc có chất gây nhiễu.
\item Phản ứng Tollens nhận biết các chất có tính khử trong điều kiện thử, không phải chỉ riêng glucose trong mọi hỗn hợp. Chỉ kết luận glucose khi biết trước phạm vi các mẫu có thể có.
\item Đông tụ khi thêm acid/nhiệt là dấu hiệu hữu ích với mẫu protein quen thuộc, nhưng không phải phép thử đặc hiệu tuyệt đối trong mọi hỗn hợp.
\item Khi làm thí nghiệm, dùng các phần mẫu riêng; thuốc thử Tollens, acid và kiềm phải do giáo viên hướng dẫn, xử lí chất thải đúng quy định.
\end{itemize}
\end{ghinho}

\subsection{Ví dụ: nhận biết một bộ mẫu thực phẩm}

Có năm mẫu riêng biệt, mỗi mẫu chỉ chứa một trong các chất sau: dầu ăn, hồ tinh bột, dung dịch glucose, dung dịch saccharose và dung dịch lòng trắng trứng. Có thể thiết kế trình tự nhận biết như sau:

\begin{enumerate}
\item \textbf{Tìm dầu ăn:} cho một ít mẫu vào nước. Mẫu tạo thành lớp riêng, không hòa tan trong nước, là ứng viên dầu ăn. Để kiểm tra bản chất triglyceride, lấy phần mẫu mới đun nóng với dung dịch \ce{NaOH}; phản ứng xà phòng hóa tạo glycerol và muối của acid béo.
\item \textbf{Tìm hồ tinh bột:} nhỏ iodine vào phần mẫu mới của các mẫu còn lại. Mẫu chuyển xanh tím là hồ tinh bột.
\item \textbf{Tìm glucose:} dùng thuốc thử Tollens với phần mẫu mới. Mẫu tạo bạc khi đun nóng là dung dịch glucose trong bộ mẫu đã cho.
\item \textbf{Phân biệt lòng trắng trứng và saccharose:} thêm một lượng nhỏ acid loãng vào hai mẫu còn lại. Mẫu protein có thể đông tụ/kết tủa; mẫu saccharose vẫn là dung dịch. Có thể dùng mẫu đối chứng và quan sát thêm khi đun nóng nhẹ theo hướng dẫn.
\end{enumerate}

Trình tự trên chỉ phù hợp với bộ mẫu đã nêu, trong đó mỗi mẫu chứa một chất. Với mẫu thực phẩm thực tế có nhiều thành phần, cần tách hoặc làm sạch mẫu trước, dùng các phần mẫu riêng và phối hợp nhiều bằng chứng thay vì kết luận từ một phép thử đơn lẻ.

\subsection{Vận dụng trong thực phẩm, nông nghiệp và sản xuất}

Trong thực phẩm, lipid cung cấp năng lượng và tạo vị béo; carbohydrate là nguồn năng lượng quan trọng; protein tham gia cấu tạo và duy trì hoạt động của cơ thể. Phân tích đúng thành phần giúp lựa chọn, chế biến và sử dụng thực phẩm hợp lí.

Trong nông nghiệp, tinh bột từ gạo, ngô, sắn và cellulose từ rơm rạ, bã mía có thể được chuyển hóa thành glucose. Glucose tiếp tục được dùng để lên men tạo ethanol hoặc tham gia phản ứng tráng bạc trong những quy trình thích hợp. Dầu, mỡ chứa triglyceride có thể được dùng làm nguyên liệu sản xuất xà phòng; protein từ đậu, sữa và phụ phẩm thực phẩm có thể được khai thác làm nguồn dinh dưỡng.

Khi phân tích một quy trình sản xuất, cần xác định đúng chất ban đầu, chất trung gian, sản phẩm, phản ứng xảy ra và hiệu suất của từng công đoạn. Không suy ra rằng mọi nguyên liệu hữu cơ đều có thể chuyển hóa bằng cùng một phản ứng.
